% GENERATED FILE. Edit canonical lessons, then run npm run generate:books.
% canonical-source-hash: fnv1a64:28d143208c7e9125
% canonical-lessons: LA-C183-notus, LA-C183-celeber, LA-C183-honestus, LA-C183-medius, LA-C183-summus

\chapter{Known, Famous, Honourable, Middle, Highest}
\label{ch:la-known-famous-honourable-middle-highest}
\hlchaptermodality{183}

\begin{chapteropening}
\textbf{By the end of this chapter:} \emph{I can say known, famous, honourable, middle and highest.}

Complete the last lesson of chapter 183: I can say known, famous, honourable, middle and highest.
\end{chapteropening}

\section[nōtus, nōta, nōtum]{nōtus, nōta, nōtum — known, famous}

\label{lesson:LA-C183-notus}



\begin{quote}
Before the new one: say the Latin for hollow, vain, then the Latin for heavy, serious.
\end{quote}

\subsection*{You'll want to know: nōtus, nōta, nōtum}
\textbf{nōtus, nōta, nōtum} — \textquotedblleft{}known, famous\textquotedblright{}.

\begin{cousinweb}
\textbf{nōtus} is the past participle of \emph{nōscō}, \textquotedblleft{}to get to know\textquotedblright{} (Lewis \& Short): known. \emph{nōscō}, older \emph{gnōscō}, is usually traced to Proto-Indo-European \emph{*gneh3-}, \textquotedblleft{}to know\textquotedblright{}.
\end{cousinweb}

\begin{grammarlens}[title={What you've built}]
One more word for this chapter.
\end{grammarlens}

\subsection*{Your turn}
\begin{itemize}
\raggedright
  \item \emph{Say it:} \emph{nōtus, nōta, nōtum}
  \item \emph{Say it:} \emph{nōtus, nōta, nōtum}, once more
  \item \emph{Say it:} say \emph{nōtus, nōta, nōtum}
  \item \emph{Recall:} say the Latin for hollow, vain, then the Latin for heavy, serious, then say \emph{nōtus, nōta, nōtum} again
\end{itemize}

\begin{tcolorbox}[breakable,colback=teal!4,colframe=teal!35!black,title={Before you move on}]
What does nōtus, nōta, nōtum mean? (\textquotedblleft{}Known, famous\textquotedblright{}.) Say it once more.
\end{tcolorbox}
\section[celeber, celebris, celebre]{celeber, celebris, celebre — famous, crowded}

\label{lesson:LA-C183-celeber}



\begin{quote}
Before the new one: say the Latin for heavy, serious, then the Latin for known.
\end{quote}

\subsection*{You'll want to know: celeber, celebris, celebre}
\textbf{celeber, celebris, celebre} — \textquotedblleft{}famous, crowded\textquotedblright{}.

\begin{cousinweb}
\textbf{celeber} has no sure origin; it first meant \textquotedblleft{}crowded, much visited\textquotedblright{}, and Lewis \& Short's comparison with a root for hearing is uncertain. English \emph{celebrate} comes from its family.
\end{cousinweb}

\begin{grammarlens}[title={What you've built}]
One more word for this chapter.
\end{grammarlens}

\subsection*{Your turn}
\begin{itemize}
\raggedright
  \item \emph{Say it:} \emph{celeber, celebris, celebre}
  \item \emph{Say it:} \emph{celeber, celebris, celebre}, once more
  \item \emph{Say it:} say \emph{celeber, celebris, celebre}
  \item \emph{Recall:} say the Latin for heavy, serious, then the Latin for known, then say \emph{celeber, celebris, celebre} again
\end{itemize}

\begin{tcolorbox}[breakable,colback=teal!4,colframe=teal!35!black,title={Before you move on}]
What does celeber, celebris, celebre mean? (\textquotedblleft{}Famous, crowded\textquotedblright{}.) Say it once more.
\end{tcolorbox}
\section[honestus, honesta, honestum]{honestus, honesta, honestum — honourable}

\label{lesson:LA-C183-honestus}



\begin{quote}
Before the new one: say the Latin for known, then the Latin for famous.
\end{quote}

\subsection*{You'll want to know: honestus, honesta, honestum}
\textbf{honestus, honesta, honestum} — \textquotedblleft{}honourable\textquotedblright{}.

\begin{cousinweb}
\textbf{honestus} is built on \emph{honōs}, \textquotedblleft{}honour, office\textquotedblright{} (Lewis \& Short), the older form of honor: endowed with honour. Where \emph{honōs} comes from is unknown.
\end{cousinweb}

\begin{grammarlens}[title={What you've built}]
One more word for this chapter.
\end{grammarlens}

\subsection*{Your turn}
\begin{itemize}
\raggedright
  \item \emph{Say it:} \emph{honestus, honesta, honestum}
  \item \emph{Say it:} \emph{honestus, honesta, honestum}, once more
  \item \emph{Say it:} say \emph{honestus, honesta, honestum}
  \item \emph{Recall:} say the Latin for known, then the Latin for famous, then say \emph{honestus, honesta, honestum} again
\end{itemize}

\begin{tcolorbox}[breakable,colback=teal!4,colframe=teal!35!black,title={Before you move on}]
What does honestus, honesta, honestum mean? (\textquotedblleft{}Honourable\textquotedblright{}.) Say it once more.
\end{tcolorbox}
\section[medius, media, medium]{medius, media, medium — middle, the middle of}

\label{lesson:LA-C183-medius}



\begin{quote}
Before the new one: say the Latin for famous, then the Latin for honourable.
\end{quote}

\subsection*{You'll want to know: medius, media, medium}
\textbf{medius, media, medium} — \textquotedblleft{}middle, the middle of\textquotedblright{}.

\begin{cousinweb}
\textbf{medius} is inherited, usually traced to Proto-Indo-European \emph{*medhyos}, \textquotedblleft{}middle\textquotedblright{} (EtymDB 2.1); Greek \emph{mesos}, \textquotedblleft{}middle\textquotedblright{}, and English \emph{mid} are cousins (Lewis \& Short).
\end{cousinweb}

\begin{grammarlens}[title={What you've built}]
One more word for this chapter.
\end{grammarlens}

\subsection*{Your turn}
\begin{itemize}
\raggedright
  \item \emph{Say it:} \emph{medius, media, medium}
  \item \emph{Say it:} \emph{medius, media, medium}, once more
  \item \emph{Say it:} say \emph{medius, media, medium}
  \item \emph{Recall:} say the Latin for famous, then the Latin for honourable, then say \emph{medius, media, medium} again
\end{itemize}

\begin{tcolorbox}[breakable,colback=teal!4,colframe=teal!35!black,title={Before you move on}]
What does medius, media, medium mean? (\textquotedblleft{}Middle, the middle of\textquotedblright{}.) Say it once more.
\end{tcolorbox}
\section[summus, summa, summum]{summus, summa, summum — highest, the top of}

\label{lesson:LA-C183-summus}



\begin{quote}
Before the new one: say the Latin for honourable, then the Latin for middle.
\end{quote}

\subsection*{You'll want to know: summus, summa, summum}
\textbf{summus, summa, summum} — \textquotedblleft{}highest, the top of\textquotedblright{}.

\begin{cousinweb}
\textbf{summus} is the superlative belonging to \emph{superus}, \textquotedblleft{}upper\textquotedblright{} (Lewis \& Short), from an older \emph{*supimos}, usually traced to Proto-Indo-European \emph{*(s)up-}, \textquotedblleft{}up, over\textquotedblright{} (EtymDB 2.1).
\end{cousinweb}

\begin{grammarlens}[title={What you've built}]
One more word for this chapter.
\end{grammarlens}

\subsection*{Your turn}
\begin{itemize}
\raggedright
  \item \emph{Say it:} \emph{summus, summa, summum}
  \item \emph{Say it:} \emph{summus, summa, summum}, once more
  \item \emph{Say it:} say \emph{summus, summa, summum}
  \item \emph{Recall:} say the Latin for honourable, then the Latin for middle, then say \emph{summus, summa, summum} again
\end{itemize}

\begin{tcolorbox}[breakable,colback=teal!4,colframe=teal!35!black,title={Before you move on}]
What does summus, summa, summum mean? (\textquotedblleft{}Highest, the top of\textquotedblright{}.) Say it once more.
\end{tcolorbox}
